\documentclass[12pt, a4paper]{article}

\usepackage{makra}

\author{M.~Godlewska}
\title{Ćwiczenia z \LaTeXe}
\date{01.12.2014}

\linespread{1.2}

\begin{document}
\maketitle

\tableofcontents

\section{Mój pierwszy rozdział}
\label{sec:MojPierwszyRozdzial}

\TeX~jest programem komputerowym stworzonym przez Donalda E. Knutha. Jest przeznaczony do składu tekstów oraz wzorów matematycznych.
Knuth rozpoczał prace nad \TeX em w 1977 roku, aby wykorzystać potencjał
składu cyfrowego, stosowanego wówczas na coraz szerszą skalę w poligrafii.
Miał też nadzieję, że uda się odwrócić tendencję do pogarszania się jakości
typograficznej, co uwidaczniało się w jego własnych książkach i artykułach.
W używanej obecnie postaci \TeX~został udostępniony w roku 1982, a niewielkie
rozszerzenie, dotyczące ośmiobitowego kodowania znaków, pojawiło
się w roku 1989. \TeX~ma renomę programu nadzwyczaj stabilnego, pracującego
na różnego rodzaju sprzęcie oraz praktycznie wolnego od błedów.
Numery wersji \TeX a zbiegają do liczby $\pi$, a obecny wynosi 3,14159.

Słowo \TeX~należy wymawiać „tech”. Zgłoska „ch” ma związek z tym, że
znak X przypomina grecką literę „chi” ($\chi$). \TeX~jest też pierwszą sylabą greckiego
słowa texnologia (technologia). W sytuacjach, w których nie można
zapisywac nazwy \TeX~z charakterystycznym obniżeniem litery E, należy
zamiennie używać wersji TeX \cite{niekrotkie}.

Teraz będzie rozdział \ref{sec:MatmaWLaTeXU}

\section{Matma w \LaTeX u}
\label{sec:MatmaWLaTeXU}

Oto kilka wzorków \eqref{eq:nr1} i \eqref{eq:nr2}:\\
najpierw $\frac{1}{8}$  i jedziemy dalej.
$$\frac{1}{9}$$
	\[\frac{1}{a^2}\]
\begin{equation}
\label{eq:nr1}
	\frac{1}{\cos{x}}
\end{equation}

\begin{equation}
\label{eq:nr2}
	\frac{1}{\sqrt{2}}
\end{equation}

\subsection{Inna matma}
\label{sec:InnaMatma}

Rozwiąż układ równań \eqref{eq:nr3}:

\begin{equation}
	y = \left\{	
	\begin{array}{ll}
	2x_1^2 + x_2 & \textrm{jeżeli $x_1 > 0 \land x_2 >0$} \\
	3x_1 - 5 x_2 & \textrm{jeżeli $x_1 < 0 \land x_2 <0$} \\
	0 & \textrm{wpp}
		
	\end{array} \right.
	\label{eq:nr3}
\end{equation}
\\
A teraz dana jest macierz:
\begin{displaymath}
\mathbf{X} =
\left[ \begin{array}{ccc}
x_{11} & x_{12} & \ldots \\
x_{21} & x_{22} & \ldots \\
\vdots & \vdots & \ddots
\end{array} \right]
\end{displaymath}

\newpage

\section{Rysunki}
\label{sec:Rysunki}
Jaki fajny jest:
\begin{figure}[h]
	\centering
		\includegraphics[width=0.50\textwidth]{images.jpg}
	\caption{Fajny \LaTeX}
	\label{fig:images}
\end{figure}
\\
A to mój własny rysunek wektorowy:
\begin{figure}[!ht]
	\centering
		\includegraphics[width=0.40\textwidth]{sloneczko.pdf}
	\caption{Słoneczko}
	\label{fig:sloneczko}
\end{figure}
\\
Po rysunkach \ref{fig:images} i \ref{fig:sloneczko} może coś zakodujemy:
\begin{figure}[!ht]
\centering
\setlength{\unitlength}{0.75mm}
\begin{picture}(60,40)
\put(30,20){\vector(1,0){30}}
\put(30,20){\vector(4,1){20}}
\put(30,20){\vector(3,1){25}}
\put(30,20){\vector(2,1){30}}
\put(30,20){\vector(1,2){10}}
\thicklines
\put(30,20){\vector(-4,1){30}}
\put(30,20){\vector(-1,4){5}}
\thinlines
\put(30,20){\vector(-1,-4){5}}
\end{picture}
\caption{Strzałeczki \cite{niekrotkie}}
\label{fig:strzaleczki}
\end{figure}

\section{Punktujemy}
\label{sec:Punktujemy}
Lista punktowana:
\begin{itemize}
	\item pierwszy
	\item[--] drugi
	\item[+] trzeci
\end{itemize}
Lista numerowana:
\begin{enumerate}
	\item coś
	
	\begin{enumerate}
		\item coś innego
		\item jeszcze coś innego
		
	\end{enumerate}
	\item i jeszcze coś 
\end{enumerate}

\bibliographystyle{plain}
\bibliography{mojaBib}
 
\end{document}
